% ============================================================
%  PLAN FOR THE REST OF THE THESIS
%  "Offensive Corner Tactics vs. Defensive Structures:
%   An Automated Classification and Effectiveness Analysis"
% ============================================================

\chapter{Plan for the Rest of the Thesis}
\label{ch:plan}

This chapter outlines the concrete next steps that will be completed
before the thesis is submitted. The work is sequenced such that the
foundational models (SQ1 and SQ2) are fully validated and applied
before moving into the inferential analyses of SQ3 and SQ4, and before
the D\&V Analyser is populated with the final results.

% ────────────────────────────────────────────────────────────────────────────
\section{Inter-rater Reliability Testing}
\label{sec:plan_interrater}

Before the manual labels used to train and evaluate the SQ1 and SQ2
models are considered final, formal inter-rater reliability testing will
be conducted. A second rater will independently label a stratified
sample of corners from the 67-corner test set, covering all major role
classes (ZONAL, MAN, TARGET, DECOY, BLOCK\_GK, BLOCK\_DEF). Agreement
will be measured using Cohen's Kappa. The purpose is twofold: to
quantify the inherent ambiguity in role assignment (providing a human
ceiling for classifier accuracy) and to identify which role categories
are most susceptible to label noise so that the model's reported
accuracy figures can be interpreted appropriately.

% ────────────────────────────────────────────────────────────────────────────
\section{Expanding the Labelled Test Set}
\label{sec:plan_labelling}

The current test set of 67 corners is sufficient for a proof-of-concept
evaluation but is small for per-class analysis, particularly for the
rare roles (BLOCK\_GK, 13 examples; BLOCK\_DEF, 14 examples). Additional
labelling of 30--50 corners will be carried out, prioritising corners
from matches that contain known block-based or highly dynamic tactics
(identified via the automated pre-classifier). Expanding the test set
will directly improve the reliability of per-class accuracy estimates
and will provide more training data for the augmentation pipeline.

% ────────────────────────────────────────────────────────────────────────────
\section{Finalising SQ1 and SQ2 Models}
\label{sec:plan_models}

\subsection*{SQ1 — Defensive setup classifier}

The SQ1 model will be evaluated against the expanded test set and any
remaining per-class weaknesses addressed. Once the accuracy meets the
acceptance threshold, the classifier will be applied to all 1{,}893
corners in the full dataset to generate defensive setup predictions at
scale.

\subsection*{SQ2 — Attacker role and strategy classifier}

The attacker role classifier (TARGET / DECOY / STATIC / SECOND\_BALL /
BLOCK\_GK / BLOCK\_DEF) is currently under active development. Two
improvements are being finalised:
\begin{enumerate}
  \item \textbf{Data augmentation} — y-coordinate flipping (corners
        are pitch-symmetric) and SMOTE over-sampling within each
        LOCO-CV fold to address the severe class imbalance in
        BLOCK\_GK and BLOCK\_DEF. Preliminary results from this scheme
        are expected to improve overall 6-class accuracy above 80\%.
  \item \textbf{Running-line integration} — a geometric running-line
        classifier assigns each attacker to a zone trajectory
        (NEAR\_POST, FAR\_POST, CENTRAL, CROSSING, etc.). These labels
        will be combined with the role predictions to populate the
        full taxonomy (Tiers C--E of the 21-category scheme), enabling
        the more granular strategic categories (e.g.\ NEAR\_POST\_TARGET,
        SCISSORS\_NF) to be assigned.
\end{enumerate}

Once both components are finalised, the model will be applied to all
1{,}401 direct corners to produce a full role-and-strategy classification
for every corner in the dataset.

% ────────────────────────────────────────────────────────────────────────────
\section{SQ3 — Effectiveness of Offensive Strategies Against Defensive Structures}
\label{sec:plan_sq3}

With the full-dataset role, strategy, and defensive-setup predictions
in hand, SQ3 can be addressed. The analysis will examine interactions
between offensive corner strategy (as classified by SQ2) and defensive
structure (as classified by SQ1) as predictors of corner effectiveness
(Corner Threat Score, xA, shot generation). Concretely:

\begin{itemize}
  \item Logistic regression on shot generation and OLS regression on
        CTS with delivery type, ball height, dominant run-line category,
        and defensive setup type as predictors.
  \item Interaction terms between offensive strategy categories (e.g.\
        BLOCK\_GK\_WAVE, MULTI\_TARGET\_RUSH) and defensive structure
        (ZONAL vs.\ MAN) to test whether specific attacking patterns are
        differentially effective against specific defences.
  \item Robustness checks: hurdle model (separate shot-generation and
        conditional-CTS components) and per-team fixed effects to
        control for opponent quality.
\end{itemize}

% ────────────────────────────────────────────────────────────────────────────
\section{SQ4 — Individual Player Quality}
\label{sec:plan_sq4}

SQ4 will analyse individual taker and attacker quality while controlling
for the tactical and structural context identified in SQ1--SQ3. The
model for SQ4 will include:

\begin{itemize}
  \item A \textbf{taker quality metric} (TakerCTS) based on the
        leave-one-match-out residual after controlling for delivery
        zone, ball height, and the opposing defensive structure. This
        isolates individual contribution from tactical context.
  \item An analysis of \textbf{attacker quality}, identifying which
        players are disproportionately often assigned the TARGET or
        BLOCK\_GK role and how their presence in the lineup affects
        corner threat.
  \item A test of whether individual quality effects survive when
        controlling for tactics and defensive structure (SQ3 covariates),
        to address the question of whether ``it's the player or the
        plan'' that drives corner effectiveness.
\end{itemize}

% ────────────────────────────────────────────────────────────────────────────
\section{Data \& Video Analyser — Final Development}
\label{sec:plan_dv}

The D\&V Analyser will be extended in two directions once the research
results are available:

\begin{enumerate}
  \item \textbf{Research results integration} — The analytical outputs
        of SQ1--SQ4 (delivery effectiveness rankings, team defensive
        profiles, attacker role predictions, strategy category per
        corner) will be surfaced inside the platform. Coaching staff
        will be able to query, for example, ``which delivery type is
        most effective when this opponent uses zonal marking?'' and
        navigate directly to video examples.
  \item \textbf{FC Den Bosch operational usage} — The platform will be
        prepared for live use by the FC Den Bosch performance analysis
        team for the remainder of the season. This includes adding their
        own match data as it becomes available, ensuring that the
        system's opponent profiling and video-retrieval components are
        robust and user-friendly for non-technical staff.
\end{enumerate}

The D\&V Analyser thus serves both as a vehicle for presenting and
validating the thesis findings and as a practical deliverable for the
club.
